\newpage

% *************************
% **     JUSTIFICACION    **
% *************************
\section{Justificación}
explicar por qué vale la pena investigarlo.
\begin{itemize}
    \item ¿Por qué es importante?
    \item ¿A quién beneficia?
    \item ¿Qué aporta teórica, práctica o metodológicamente?
\end{itemize}

% *************************
% **     PROBLEMA        **
% *************************

\section{Problema de Investigación}
\subsection{Descripción del problema}
En esta sección  de detalla desde el punto de vista más genérico (mundial) hasta llegar al problema local que nos afecta.\

Nos podemos apoyar de herramientas como árbol de problemas (Efectos, Problema, Causas); o también a AQP (causas; [A donde, Quien, Problema]; Consecuencias; aporte {P y [causa - consecuencia -aporte ]en Q de A}) como posible titulo.\

\subsection{Identificación del problema}
Focalizamos el lugar de estudio como afecta y como es el comportamiento.\

\subsection{Formulación del problema}
\subsubsection{Problema general.}
La formulación del problema describe de manera clara la situación problemática que da origen a la investigación, identificando sus principales causas y consecuencias dentro de un contexto determinado. Asimismo, delimita los aspectos clave que serán abordados y justifica la necesidad de aplicar una solución.

Finalmente, el apartado culmina con la pregunta de investigación que orienta el desarrollo del estudio e incorpora la solución técnica propuesta. (usar palabras como: eficiencia, tiempo, desempeño, optimización, productividad, nos dan como respuestas soluciones tecnológicas).

\subsubsection{Sub problemas.}
Se desglosa el problema principal.


% *************************
% **     ALCANCES        **
% *************************
\section{Alcances y limitaciones}
\subsection{Alcances}
Indicar hasta donde nos proponemos llegar con el trabajo de investigación, comparación de algoritmos, mejora de los mismo, implementación de un sistema o creación de un prototipo, lo que se requiera.
\subsection{Limitaciones}
Enmarcar detalladamente el trabajo de investigación a nivel temporal y espacial.


% *************************
% **     OBJETIVOS    **
% *************************

\section{Objetivos}
\subsection{Objetivo General}
% solo uno
Lleva estricta relación con el titulo del trabajo, y responde como solución al problema identificado.

\subsection{Objetivos Específicos}
\begin{itemize}
	\item Responde como solucionar el problema especifico X.
	\item Responde como solucionar el problema especifico Y.
	\item Responde como solucionar el problema especifico Z.
\end{itemize}
